\begin{frame}{왜 ``곱''인가 — 할인율 비유}
  \begin{itemize}
    \item $L$이 층 1의 $W_1$에 영향을 미치는 경로:
      $W_1 \to z_1 \to a_1 \to z_2 \to \cdots \to L$ —
      연쇄법칙이 이 사슬을 따라 미분을 \textbf{연속으로 곱}
    \item 곱 연산은 \textbf{할인율}처럼: 각 층이 신호의 일부만 전달하면
      10층에 $0.25^{10}$으로 지수 감소
    \item 비유: 10명이 순서대로 ``들린 것의 25\%만'' 다음 사람에게
      속삭이면 — 10번째에게는 100만 분의 1도 안 남음
  \end{itemize}
  \begin{block}{핵심}
    소실은 ``느리다''가 아니라 \kb{``지수적으로 0에 수렴한다''}.
  \end{block}
  \vskip0.3em
  {\footnotesize 소실과 폭발은 \textbf{같은 연쇄법칙 곱의 두 얼굴}:
  각 층의 ``승수''가 1보다 작으면 소실, 크면 폭발. 이 승수는
  활성함수($\sigma'$)와 가중치($W_l$)의 크기로 — 이번 절의
  ``활성함수''와 ``초기화·정규화''가 바로 이 승수를 제어하는 것.}
\end{frame}
